% ===========================
% Persian Exam Template
% ===========================
\documentclass[12pt]{exam}

% ---------------------------
% 1. Packages
% ---------------------------
\usepackage{tikz}
\usetikzlibrary{arrows.meta, positioning}

\usepackage{fontspec}
\usepackage{setspace}
\usepackage{listings}
\usepackage{graphicx}
\usepackage{wrapfig}
\usepackage{subfigure}
\usepackage{multirow}
\usepackage{amsmath}
\usepackage{amssymb}
\usepackage{multicol}
\usepackage{changepage}
\usepackage[margin=20mm]{geometry}
\usepackage{physics}
\usepackage{adjustbox}
\usepackage{tabularx}
\usepackage{xcolor}
\usepackage{enumitem}

% xepersian باید آخر لود شود
\usepackage{xepersian}

% ---------------------------
% 2. Persian Font
% ---------------------------
\settextfont[
    Path=fonts/,
    Extension=.ttf
]{XB Niloofar}

% ---------------------------
% 3. Fix parts environment
% ---------------------------
% این بخش باعث می‌شود \begin{parts} و \part از exam استفاده نکنند
% و با XePersian قاطی نکنند.
\renewenvironment{parts}
{
    \begin{enumerate}[
        label=\textbf{\alph*)},
        itemsep=1em,
        topsep=0.5em,
        rightmargin=0pt,
        leftmargin=0pt
    ]
    \renewcommand{\part}{\item}
}
{
    \end{enumerate}
}

% ---------------------------
% 4. Commands
% ---------------------------
\newcommand{\bin}[1]{\texttt{#1}}

\newcommand{\namee}{کسری شیخی}
\newcommand{\idd}{۴۰۳۱۱۰۵۹۶}
\newcommand{\coursename}{هوش مصنوعی}
\newcommand{\term}{بهار ۱۴۰۵}
\newcommand{\college}{دانشکده مهندسی کامپیوتر}
\newcommand{\prof}{استاد: آرش ماری‌اوریاد}

\singlespacing
\parindent 0ex

% ---------------------------
% 5. Listings
% ---------------------------
\lstset{
    keywordstyle=\textbf,
    identifierstyle=,
    stringstyle=\ttfamily,
    commentstyle=\color{LimeGreen},
    numberstyle=\footnotesize,
    showstringspaces=false
}

\begin{document}

% ---------------------------
% Header
% ---------------------------
\pagestyle{head}
\firstpageheader{}{}{}
\runningheader{صفحه \thepage\ از \numpages}{}{\coursename}
\runningheadrule

\begin{tabular}{p{.7\textwidth} l}
\multicolumn{2}{c}{\textbf{به نام خدا}}\\
\multirow{2}{*}{\includegraphics[scale=0.2]{images/logo.png}} & \\ \\ \\ \\
& \textbf{\namee}\\
& \textbf{\idd}\\
\\
\textbf{\coursename}\\
\textbf{\term}\\
\textbf{\prof}\\ \\
\textbf{\college} &  \\
\end{tabular}\\

\vspace{5mm}
\rule[1ex]{\textwidth}{.1pt}

\textbf{پاسخ تمرین سری پنجم - یادگیری تقویتی}

\rule[1ex]{\textwidth}{.1pt}

% ---------------------------
% Questions
% ---------------------------
\section*{سؤال اول}
\begin{parts}

\part درست

 همگرایی سیاست تنها به \lr{argmax} مقادیر $Q$ در هر حالت بستگی دارد،
نه به مقدار دقیق آن‌ها. بنابراین ممکن است پس از چند ایتریشن، \lr{argmax} در همه‌ی حالت‌ها
ثابت شود و سیاست حریصانه دیگر تغییر نکند، در حالی که $V$ هنوز به مقدار دقیق $V^*$ نرسیده
و به‌روزرسانی آن ادامه دارد.

\part نادرست

 \lr{Q-learning} یک الگوریتم \lr{off-policy} است؛ یعنی مقادیر بهینه‌ی $Q$ را می‌تواند
مستقل از اینکه کنش‌های واقعاً انتخاب‌شده طبق سیاست رفتاری بهینه باشند یا نه، یاد بگیرد.
شرط کافی همگرایی، بازدید کافی از همه‌ی زوج‌های $(s,a)$ است، نه بهینه بودن سیاست رفتاری.

\part درست 

 در \lr{Value Iteration} هر ایتریشن شامل به‌روزرسانی همه‌ی حالت‌ها با پیچیدگی
$O(|S|^2|A|)$ است. در \lr{Policy Iteration} هر ایتریشن شامل ارزیابی سیاست با پیچیدگی
$O(|S|^3)$ و بهبود سیاست با پیچیدگی $O(|S|^2|A|)$ است.
اگر $|A|\approx |S|$، هر دو از مرتبه‌ی $O(|S|^3)$ می‌شوند، پس پیچیدگی هر ایتریشن تقریباً
برابر خواهد بود.

\part نادرست

 اگر به همه‌ی پاداش‌ها یک ثابت $c$ اضافه کنیم، سیاست بهینه لزوماً ثابت نمی‌ماند.
دلیلش این است که اگر \lr{MDP} حالت \lr{terminal} داشته باشد، سیاست‌هایی که دیرتر به
حالت پایانی می‌رسند تعداد پاداش‌های بیشتری می‌گیرند؛ پس مقدار اضافه‌شده برای همه‌ی سیاست‌ها
یکسان نیست.

البته اگر افق نامحدود باشد و پایان اپیزود در کار نباشد، در این حالت مقدار اضافه‌شده برای
همه‌ی سیاست‌ها برابر با
\[
\frac{c}{1-\gamma}
\]
می‌شود و آنگاه گزاره درست خواهد بود. اما در حالت کلی گزاره نادرست است.

\part نادرست

 مقدار $\gamma$ میزان اهمیت پاداش‌های آینده نسبت به پاداش‌های نزدیک را مشخص می‌کند.
پس اگر $\gamma$ عوض شود، ممکن است کنش بهینه هم عوض شود.

برای مثال یک \lr{MDP} را در نظر بگیرید که از حالت شروع $s$ دو کنش داریم:

\[
a_1:\quad \text{دریافت پاداش } 5 \text{ و پایان}
\]

\[
a_2:\quad \text{ابتدا پاداش } 0 \text{ و سپس در گام بعدی پاداش } 10 \text{ و پایان}
\]

ارزش کنش اول برابر است با:

\[
Q(s,a_1)=5
\]

و ارزش کنش دوم برابر است با:

\[
Q(s,a_2)=10\gamma
\]

اگر $\gamma=0.4$ باشد، داریم:

\[
Q(s,a_2)=4<5
\]

پس کنش $a_1$ بهتر است. اما اگر $\gamma=0.9$ باشد، داریم:

\[
Q(s,a_2)=9>5
\]

پس کنش $a_2$ بهتر می‌شود. بنابراین دو \lr{MDP} که فقط در مقدار $\gamma$ فرق دارند،
لزوماً سیاست بهینه‌ی یکسانی ندارند.

\part درست

 در افق محدود (\lr{finite horizon})، ارزش هر کنش به تعداد گام‌های باقی‌مانده تا پایان
افق بستگی دارد. بنابراین حتی برای یک حالت ثابت $s$، سیاست بهینه بسته به این‌که چند گام تا
پایان مانده می‌تواند فرق کند؛ یعنی سیاست بهینه \lr{non-stationary} است، برخلاف \lr{MDP}
با افق نامحدود که سیاست بهینه \lr{stationary} است.
\end{parts}


\vspace{1cm}

\begin{parts}

\part
تفاوت \lr{on-policy RL} و \lr{off-policy RL}:
در روش‌های \lr{on-policy}، عامل ارزش سیاستی را یاد می‌گیرد که برای تولید تجربه از همان
استفاده می‌کند؛ یعنی سیاست رفتاری (\lr{behavior policy}) و سیاست هدف (\lr{target policy})
یکسان هستند، مانند \lr{SARSA}. در روش‌های \lr{off-policy}، سیاست تولید تجربه می‌تواند با
سیاستی که عامل قصد یادگیری ارزش آن را دارد متفاوت باشد، مانند \lr{Q-learning} که $Q^*$ را
یاد می‌گیرد، حتی اگر کنش‌ها با سیاستی مثل \lr{$\varepsilon$-greedy} انتخاب شده باشند.

\part
تفاوت \lr{model-based RL} و \lr{model-free RL}:
در \lr{model-based RL}، عامل مدل محیط یعنی توابع گذار $T(s'\mid s,a)$ و پاداش
$R(s,a,s')$ را از قبل می‌داند یا از تجربه تخمین می‌زند و بعد با استفاده از این مدل
برنامه‌ریزی می‌کند. در \lr{model-free RL}، عامل بدون ساختن مدل صریحی از محیط، مستقیماً
از تعامل با محیط تابع ارزش یا سیاست را یاد می‌گیرد، مانند \lr{Q-learning}، \lr{SARSA}
و \lr{Monte Carlo}.

\end{parts}
\clearpage

\section*{سؤال دوم}
\begin{parts}

\part
رابطه به‌روزرسانی در \lr{Value Iteration} برابر است با:
\[
V_{k+1}(s)=
\max_a
\sum_{s'}T(s'|s,a)\left(R(s,a,s')+\gamma V_k(s')\right),
\qquad
\gamma=0.9
\]
در ابتدا:
\[
V_0(s)=0
\]
حالت‌های \lr{terminal} مقدار صفر دارند:
\[
V(3,1)=0,\qquad V(3,2)=0
\]

\textbf{مرحله اول:}

فقط دو حالت می‌توانند در یک حرکت به خانه با پاداش مثبت برسند:
\[
V_1(2,2)=0.9(10)=9,\qquad
V_1(3,3)=0.9(10)=9
\]
پس:
\[
V_1(1,1)=0,\quad
V_1(1,2)=0,\quad
V_1(1,3)=0,\quad
V_1(2,1)=0
\]
\[
V_1(2,2)=9,\quad
V_1(2,3)=0,\quad
V_1(3,3)=9
\]

\textbf{مرحله دوم:}
\[
V_2(1,1)=0
\]
\[
V_2(1,2)=0.9\left(0.9V_1(2,2)\right)=0.9(0.9\times 9)=7.29
\]
\[
V_2(1,3)=0
\]
\[
V_2(2,1)=0.9\left(0.9V_1(2,2)\right)+0.05(-10)
=7.29-0.5=6.79
\]
\[
V_2(2,2)=0.9(10)=9
\]
\[
V_2(2,3)
=
0.9\left(0.9V_1(3,3)\right)
+
0.05\left(0.9V_1(2,2)\right)
=
7.29+0.405=7.695
\]
\[
V_2(3,3)
=
0.9(10)+0.05\left(0.9V_1(3,3)\right)
=
9+0.405=9.405
\]

پس:
\[
V_2(1,1)=0,\quad
V_2(1,2)=7.29,\quad
V_2(1,3)=0
\]
\[
V_2(2,1)=6.79,\quad
V_2(2,2)=9,\quad
V_2(2,3)=7.695,\quad
V_2(3,3)=9.405
\]

\part
در \lr{TD Learning} داریم:
\[
V(s)\leftarrow V(s)+\alpha\left(r+\gamma V(s')-V(s)\right),
\qquad
\alpha=0.1,\quad \gamma=0.9
\]
در ابتدا همه مقدارها صفر هستند.

\textbf{مسیر اول:}
\[
(1,1)\to(1,2)\to(1,3)
\]
در این مسیر \lr{reward}ای دریافت نمی‌شود:
\[
V(1,1)\leftarrow 0+0.1\left(0+0.9V(1,2)-0\right)=0
\]
\[
V(1,2)\leftarrow 0+0.1\left(0+0.9V(1,3)-0\right)=0
\]

\textbf{مسیر دوم:}
\[
(1,1)\to(1,2)\to(2,2)\to(2,3)
\]
باز هم \lr{reward}ای دریافت نمی‌شود:
\[
V(1,1)\leftarrow 0+0.1\left(0+0.9V(1,2)-0\right)=0
\]
\[
V(1,2)\leftarrow 0+0.1\left(0+0.9V(2,2)-0\right)=0
\]
\[
V(2,2)\leftarrow 0+0.1\left(0+0.9V(2,3)-0\right)=0
\]

پس:
\[
V(1,1)=0,\qquad V(1,2)=0,\qquad V(2,2)=0
\]
و بقیه مقدارها نیز صفر می‌مانند.

\part
در \lr{SARSA} داریم:
\[
Q(s,a)\leftarrow Q(s,a)+\alpha\left(r+\gamma Q(s',a')-Q(s,a)\right),
\qquad
\alpha=0.1,\quad \gamma=0.9
\]
در ابتدا:
\[
Q_0(s,a)=0
\]
کنش راست را با \(a_R\) نشان می‌دهیم.

\textbf{مسیر اول:}
\[
(1,1)\to(1,2)\to(1,3)
\]
\[
Q((1,1),a_R)
\leftarrow
0+0.1\left(0+0.9Q((1,2),a_R)-0\right)=0
\]
\[
Q((1,2),a_R)
\leftarrow
0+0.1\left(0+0.9Q((1,3),a_R)-0\right)=0
\]

\textbf{مسیر دوم:}
\[
(1,1)\to(1,2)\to(2,2)\to(2,3)
\]
\[
Q((1,1),a_R)
\leftarrow
0+0.1\left(0+0.9Q((1,2),a_R)-0\right)=0
\]
\[
Q((1,2),a_R)
\leftarrow
0+0.1\left(0+0.9Q((2,2),a_R)-0\right)=0
\]
\[
Q((2,2),a_R)
\leftarrow
0+0.1\left(0+0.9Q((2,3),a_R)-0\right)=0
\]

پس:
\[
Q((1,1),a_R)=0,\qquad
Q((1,2),a_R)=0,\qquad
Q((2,2),a_R)=0
\]

\part
در \lr{Q-learning} داریم:
\[
Q(s,a)\leftarrow
Q(s,a)+\alpha\left(r+\gamma\max_{a'}Q(s',a')-Q(s,a)\right)
\]
تفاوتش با \lr{SARSA} این است که در \lr{SARSA} از \lr{action} واقعی بعدی استفاده می‌شود، ولی در \lr{Q-learning} بهترین \lr{action} ممکن در حالت بعدی را می‌گیریم.

چون مقدار اولیه همه \(Q\)ها صفر است و در این دو مسیر \lr{reward} غیرصفری نداریم، همه آپدیت‌ها صفر می‌شوند:
\[
Q((1,1),a_R)=0,\qquad
Q((1,2),a_R)=0,\qquad
Q((2,2),a_R)=0
\]

پس در این مثال مقدار عددی \lr{SARSA} و \lr{Q-learning} یکی می‌شود، ولی از نظر مفهومی یکی نیستند.

\part
در \lr{Direct Evaluation} برای هر حالتی که در مسیر دیده شده، \lr{return} تخفیف‌خورده را حساب می‌کنیم:
\[
G_t=\sum_{k=0}^{T-t-1}\gamma^k r_{t+k+1},
\qquad
\gamma=0.9
\]

\textbf{مسیر اول:}
\[
(1,1)\to(1,2)\to(1,3)
\]
چون \lr{reward} نداریم:
\[
G(1,1)=0,\qquad G(1,2)=0
\]

\textbf{مسیر دوم:}
\[
(1,1)\to(1,2)\to(2,2)\to(2,3)
\]
باز هم \lr{reward} نداریم:
\[
G(1,1)=0,\qquad G(1,2)=0,\qquad G(2,2)=0
\]

پس:
\[
\hat{V}(1,1)=\frac{0+0}{2}=0
\]
\[
\hat{V}(1,2)=\frac{0+0}{2}=0
\]
\[
\hat{V}(2,2)=0
\]

حالت‌های \((1,3)\) و \((2,3)\) انتهای مسیرهای مشاهده‌شده هستند، پس \lr{return} قابل‌استفاده‌ای برایشان نداریم. حالت‌هایی مثل \((2,1)\) و \((3,3)\) هم دیده نشده‌اند.

از نظر \lr{sample efficiency}، روش \lr{Direct Evaluation} ضعیف‌تر از \lr{TD Learning} است، چون باید ادامه مسیر را ببیند. ولی \lr{TD Learning} با هر گام می‌تواند آپدیت انجام دهد و اطلاعات را سریع‌تر پخش کند.

\end{parts}
\clearpage

\section*{سؤال سوم}
\begin{parts}

\part
می‌خواهیم بررسی کنیم آیا برای سیاست \(\pi_3\) داریم:

\[
V_3 = \frac{1}{2}V_1 + \frac{1}{2}V_2
\]

پاسخ منفی است. برای رد این گزاره یک مثال نقض می‌زنیم.

یک \lr{MDP} با یک حالت غیرپایانی \(s\) و یک حالت پایانی \(t\) در نظر می‌گیریم.
در حالت \(s\) دو کنش \(a\) و \(b\) داریم:

\[
a: \quad s \to s,\qquad r=0
\]

\[
b: \quad s \to t,\qquad r=1
\]

حالت \(t\) پایانی است و ارزش آن صفر است. ضریب تخفیف را هم برابر
\(\gamma = 0.9\) در نظر می‌گیریم.

سیاست اول \(\pi_1\) همیشه کنش \(a\) را انتخاب می‌کند. پس عامل همیشه در حالت
\(s\) می‌ماند و هیچ پاداشی نمی‌گیرد. بنابراین:

\[
V_1(s)=0
\]

سیاست دوم \(\pi_2\) همیشه کنش \(b\) را انتخاب می‌کند. پس عامل همان اول پاداش
\(1\) می‌گیرد و وارد حالت پایانی می‌شود. بنابراین:

\[
V_2(s)=1
\]

اگر رابطه‌ی داده‌شده درست بود، باید می‌داشتیم:

\[
\frac{1}{2}V_1(s)+\frac{1}{2}V_2(s)
=
\frac{1}{2}(0)+\frac{1}{2}(1)
=
\frac{1}{2}
\]

حالا سیاست \(\pi_3\) را بررسی می‌کنیم. این سیاست در هر بار حضور در حالت \(s\)،
با احتمال \(\frac{1}{2}\) کنش \(a\) و با احتمال \(\frac{1}{2}\) کنش \(b\) را انتخاب می‌کند.

پس معادله بلمن برای \(V_3(s)\) به صورت زیر است:

\[
V_3(s)
=
\frac{1}{2}\bigl(0+\gamma V_3(s)\bigr)
+
\frac{1}{2}\bigl(1+\gamma V_3(t)\bigr)
\]

چون \(t\) حالت پایانی است، داریم:

\[
V_3(t)=0
\]

پس:

\[
V_3(s)
=
\frac{1}{2}\gamma V_3(s)
+
\frac{1}{2}
\]

با جایگذاری \(\gamma=0.9\):

\[
V_3(s)
=
0.45V_3(s)+0.5
\]

\[
\begin{aligned}
V_3(s)-0.45V_3(s) &= 0.5 \\
0.55V_3(s) &= 0.5 \\
V_3(s) &= \frac{0.5}{0.55} \\
V_3(s) &= \frac{10}{11}
\end{aligned}
\]

اما از طرف دیگر داشتیم:

\[
\frac{1}{2}V_1(s)+\frac{1}{2}V_2(s)=\frac{1}{2}
\]

پس:

\[
V_3(s)=\frac{10}{11}
\neq
\frac{1}{2}
=
\frac{1}{2}V_1(s)+\frac{1}{2}V_2(s)
\]

بنابراین نمی‌توان گفت تابع ارزش سیاست مخلوط، حتماً برابر با میانگین تابع ارزش
دو سیاست اولیه است. دلیلش هم این است که سیاست جدید در هر گام دوباره تصمیم
تصادفی می‌گیرد و این باعث می‌شود مسیرهایی ساخته شوند که رفتارشان صرفاً میانگین
دو سیاست قبلی نیست.

\end{parts}
\clearpage

\section*{سؤال چهارم}
\begin{enumerate}

\item[\textbf{آ)}]

در این سؤال ضریب تخفیف برابر است با:

\[
\gamma = 0.5
\]

سیاست اولیه \(\pi_1\) در هر حالت، بین دو کنش \(x\) و \(y\) با احتمال برابر انتخاب می‌کند. همچنین مقدار اولیه‌ی تابع ارزش برای همه‌ی حالت‌ها برابر است با:

\[
V_1(A)=V_1(B)=V_1(C)=2
\]

برای اجرای یک مرحله از ارزیابی سیاست، مقدار هر حالت را با استفاده از مقدارهای قبلی حساب می‌کنیم.

\medskip

\textbf{حالت \(A\):}

اگر کنش \(x\) را انتخاب کنیم، با پاداش \(4\) به حالت \(B\) می‌رویم. پس:

\[
Q_1(A,x)=4+\gamma V_1(B)=4+0.5(2)=5
\]

اگر کنش \(y\) را انتخاب کنیم، با پاداش \(2\) به حالت \(C\) می‌رویم. پس:

\[
Q_1(A,y)=2+\gamma V_1(C)=2+0.5(2)=3
\]

چون سیاست تصادفی است، مقدار جدید \(A\) برابر میانگین این دو مقدار می‌شود:

\[
V_2(A)=\frac{1}{2}Q_1(A,x)+\frac{1}{2}Q_1(A,y)
=
\frac{1}{2}(5)+\frac{1}{2}(3)=4
\]

\medskip

\textbf{حالت \(B\):}

اگر کنش \(x\) را انتخاب کنیم، با پاداش \(2\) به حالت \(C\) می‌رویم:

\[
Q_1(B,x)=2+\gamma V_1(C)=2+0.5(2)=3
\]

اگر کنش \(y\) را انتخاب کنیم، با پاداش \(0\) به حالت \(A\) می‌رویم:

\[
Q_1(B,y)=0+\gamma V_1(A)=0+0.5(2)=1
\]

پس:

\[
V_2(B)=\frac{1}{2}Q_1(B,x)+\frac{1}{2}Q_1(B,y)
=
\frac{1}{2}(3)+\frac{1}{2}(1)=2
\]

\medskip

\textbf{حالت \(C\):}

اگر کنش \(y\) را انتخاب کنیم، با پاداش \(1\) به حالت \(B\) می‌رویم:

\[
Q_1(C,y)=1+\gamma V_1(B)=1+0.5(2)=2
\]

اگر کنش \(x\) را انتخاب کنیم، پاداش \(8\) می‌گیریم و سپس با احتمال \(\frac{1}{4}\) به \(A\) می‌رویم و با احتمال \(\frac{3}{4}\) در \(C\) می‌مانیم. بنابراین:

\[
Q_1(C,x)
=
8+\gamma\left(\frac{1}{4}V_1(A)+\frac{3}{4}V_1(C)\right)
\]

\[
Q_1(C,x)
=
8+0.5\left(\frac{1}{4}(2)+\frac{3}{4}(2)\right)
=
8+0.5(2)=9
\]

پس:

\[
V_2(C)=\frac{1}{2}Q_1(C,x)+\frac{1}{2}Q_1(C,y)
=
\frac{1}{2}(9)+\frac{1}{2}(2)=5.5
\]

در نتیجه بعد از یک مرحله ارزیابی سیاست داریم:

\[
\boxed{
V_2(A)=4,\qquad V_2(B)=2,\qquad V_2(C)=5.5
}
\]

\item[\textbf{ب)}]

حالا باید نسبت به تابع ارزش جدید، یعنی \(V_2\)، سیاست حریصانه‌ی جدید \(\pi_2\) را پیدا کنیم.

\[
V_2(A)=4,\qquad V_2(B)=2,\qquad V_2(C)=5.5
\]

\medskip

\textbf{حالت \(A\):}

\[
Q_2(A,x)=4+\gamma V_2(B)=4+0.5(2)=5
\]

\[
Q_2(A,y)=2+\gamma V_2(C)=2+0.5(5.5)=4.75
\]

چون:

\[
Q_2(A,x)>Q_2(A,y)
\]

پس:

\[
\pi_2(A)=x
\]

\medskip

\textbf{حالت \(B\):}

\[
Q_2(B,x)=2+\gamma V_2(C)=2+0.5(5.5)=4.75
\]

\[
Q_2(B,y)=0+\gamma V_2(A)=0+0.5(4)=2
\]

چون:

\[
Q_2(B,x)>Q_2(B,y)
\]

پس:

\[
\pi_2(B)=x
\]

\medskip

\textbf{حالت \(C\):}

\[
Q_2(C,y)=1+\gamma V_2(B)=1+0.5(2)=2
\]

\[
Q_2(C,x)
=
8+\gamma\left(\frac{1}{4}V_2(A)+\frac{3}{4}V_2(C)\right)
\]

\[
Q_2(C,x)
=
8+0.5\left(\frac{1}{4}(4)+\frac{3}{4}(5.5)\right)
\]

\[
Q_2(C,x)
=
8+0.5(1+4.125)
=
8+2.5625
=
10.5625
\]

چون:

\[
Q_2(C,x)>Q_2(C,y)
\]

پس:

\[
\pi_2(C)=x
\]

در نتیجه سیاست حریصانه‌ی جدید برابر است با:

\[
\boxed{
\pi_2(A)=x,\qquad \pi_2(B)=x,\qquad \pi_2(C)=x
}
\]

\item[\textbf{ج)}]

می‌خواهیم نشان دهیم اگر سیاست \(\pi'\) به صورت حریصانه از روی \(V^\pi\) ساخته شده باشد، آنگاه برای همه‌ی حالت‌ها داریم:

\[
V^{\pi'}(s)\ge V^\pi(s)
\]

چون \(\pi'\) نسبت به \(V^\pi\) حریصانه است، برای هر حالت \(s\) داریم:

\[
\pi'(s)=\arg\max_a Q^\pi(s,a)
\]

پس:

\[
Q^\pi(s,\pi'(s))=\max_a Q^\pi(s,a)
\]

از طرفی کنشی که سیاست \(\pi\) انتخاب می‌کند، یکی از همین کنش‌هاست. بنابراین:

\[
Q^\pi(s,\pi'(s))\ge Q^\pi(s,\pi(s))
\]

ولی طبق تعریف تابع ارزش سیاست \(\pi\) داریم:

\[
Q^\pi(s,\pi(s))=V^\pi(s)
\]

پس نتیجه می‌شود:

\[
Q^\pi(s,\pi'(s))\ge V^\pi(s)
\]

حالا \(Q^\pi(s,\pi'(s))\) را با رابطه‌ی بلمن باز می‌کنیم:

\[
Q^\pi(s,\pi'(s))
=
R(s,\pi'(s))
+
\gamma
\sum_{s'}P(s'|s,\pi'(s))V^\pi(s')
\]

پس داریم:

\[
V^\pi(s)
\le
R(s,\pi'(s))
+
\gamma
\sum_{s'}P(s'|s,\pi'(s))V^\pi(s')
\]

حالا اگر همین نامساوی را برای حالت‌های بعدی هم به صورت تکراری جایگذاری کنیم، سمت راست به ارزش اجرای سیاست \(\pi'\) تبدیل می‌شود. یعنی:

\[
V^\pi(s)
\le
\mathbb{E}_{\pi'}\left[
\sum_{t=0}^{\infty}\gamma^t R_t
\mid S_0=s
\right]
\]

پس:

\[
V^\pi(s)\le V^{\pi'}(s)
\]

بنابراین:

\[
\boxed{
V^{\pi'}(s)\ge V^\pi(s)
}
\]

برای بخش دوم، فرض کنیم برای همه‌ی حالت‌ها تساوی برقرار باشد:

\[
V^{\pi'}(s)=V^\pi(s)
\]

از طرفی چون \(\pi'\) نسبت به \(V^\pi\) حریصانه ساخته شده است، داریم:

\[
T^*V^\pi = T^{\pi'}V^\pi
\]

و چون طبق فرض \(V^{\pi'}=V^\pi\)، پس:

\[
V^\pi = T^{\pi'}V^\pi
\]

بنابراین:

\[
V^\pi = T^*V^\pi
\]

این یعنی \(V^\pi\) معادله‌ی بهینگی بلمن را ارضا می‌کند. از آن‌جا که جواب معادله‌ی بهینگی بلمن یکتاست، داریم:

\[
V^\pi = V^*
\]

پس سیاست \(\pi\) بهینه است. همچنین چون \(\pi'\) نسبت به \(V^*\) حریصانه است، \(\pi'\) هم سیاست بهینه خواهد بود.

در نتیجه اگر تساوی برای همه‌ی حالت‌ها رخ دهد، سیاست به دست آمده حتماً بهینه است.

\[
\boxed{
V^{\pi'}=V^\pi \Longrightarrow V^\pi=V^*
}
\]

\end{enumerate}
\clearpage

\section*{سؤال پنجم}
\begin{parts}

\part[\textbf{الف)}]
چون کارت‌ها فقط می‌توانند مقدارهای \(2\)، \(3\) و \(4\) داشته باشند، از حالت شروع فقط جمع‌های \(0,2,3,4,5\) قابل رسیدن‌اند. پس فضای حالت برابر است با:

\[
S=\{(0,C),(2,C),(3,C),(4,C),(5,C),\mathrm{Done}\}
\]

که \(C\) یعنی بازی هنوز در جریان است و \(\mathrm{Done}\) یعنی بازی تمام شده است.

\part[\textbf{ب)}]
اگر عامل کنش \lr{stop} را در حالت \((n,C)\) انتخاب کند، بازی تمام می‌شود و پاداش برابر \(n\) است:

\[
R((n,C),\mathrm{stop})=n
\]

اگر کنش \lr{draw} را انتخاب کند، پاداش آنی صفر است. اگر جمع جدید کمتر از \(6\) باشد، بازی ادامه پیدا می‌کند؛ وگرنه به حالت \(\mathrm{Done}\) می‌رود و امتیاز نهایی صفر می‌شود:

\[
R((n,C),\mathrm{draw})=0
\]

\[
R(\mathrm{Done},\cdot)=0
\]

\part[\textbf{ج)}]
برای مقدار بهینه داریم:

\[
V^*(s)=\max_a Q^*(s,a)
\]

و چون \(\gamma=1\)، محاسبات از حالت‌های آخر ساده‌تر انجام می‌شود.

برای حالت \((5,C)\):

\[
V^*(5,C)=\max\{5,0\}=5
\]

برای حالت \((4,C)\):

\[
V^*(4,C)=\max\{4,0\}=4
\]

برای حالت \((3,C)\):

\[
Q^*((3,C),\mathrm{draw})=\frac{1}{3}V^*(5,C)+\frac{1}{3}(0)+\frac{1}{3}(0)=\frac{5}{3}
\]

\[
V^*(3,C)=\max\left\{3,\frac{5}{3}\right\}=3
\]

برای حالت \((2,C)\):

\[
Q^*((2,C),\mathrm{draw})
=
\frac{1}{3}V^*(4,C)+\frac{1}{3}V^*(5,C)+\frac{1}{3}(0)
=
\frac{4+5}{3}=3
\]

\[
V^*(2,C)=\max\{2,3\}=3
\]

برای حالت \((0,C)\):

\[
Q^*((0,C),\mathrm{draw})
=
\frac{1}{3}V^*(2,C)+\frac{1}{3}V^*(3,C)+\frac{1}{3}V^*(4,C)
=
\frac{3+3+4}{3}
=
\frac{10}{3}
\]

\[
V^*(0,C)=\frac{10}{3}
\]

پس مقدارهای نهایی و سیاست بهینه به صورت زیر است:

\[
V^*(0,C)=\frac{10}{3},\qquad \pi^*(0,C)=\mathrm{draw}
\]

\[
V^*(2,C)=3,\qquad \pi^*(2,C)=\mathrm{draw}
\]

\[
V^*(3,C)=3,\qquad \pi^*(3,C)=\mathrm{stop}
\]

\[
V^*(4,C)=4,\qquad \pi^*(4,C)=\mathrm{stop}
\]

\[
V^*(5,C)=5,\qquad \pi^*(5,C)=\mathrm{stop}
\]

\[
V^*(\mathrm{Done})=0
\]

\part[\textbf{د)}]
در این بخش فرض می‌کنیم در همه‌ی نمونه‌ها، هر بار که کنش \lr{draw} زده شده، کارت \(2\) آمده است. پس از دید داده‌ها، \lr{draw} به شکل زیر عمل می‌کند:

\[
(n,C)\xrightarrow{\mathrm{draw}}(n+2,C)
\]

البته اگر \(n+2\ge 6\)، بازی تمام می‌شود. با تعداد خیلی زیاد نمونه، مقدارهای \(Q\) به مقدارهای زیر همگرا می‌شوند:

\[
Q((0,C),\mathrm{stop})=0,\qquad Q((0,C),\mathrm{draw})=4
\]

\[
Q((2,C),\mathrm{stop})=2,\qquad Q((2,C),\mathrm{draw})=4
\]

\[
Q((3,C),\mathrm{stop})=3,\qquad Q((3,C),\mathrm{draw})=5
\]

\[
Q((4,C),\mathrm{stop})=4,\qquad Q((4,C),\mathrm{draw})=0
\]

\[
Q((5,C),\mathrm{stop})=5,\qquad Q((5,C),\mathrm{draw})=0
\]

این مقدارها لزوماً برابر \(Q^*\) واقعی نیستند، چون نمونه‌ها نماینده‌ی واقعی محیط نیستند و کارت‌های \(3\) و \(4\) اصلاً دیده نشده‌اند.

\part[\textbf{ه)}]
در حالت \((4,C)\) داریم:

\[
Q((4,C),\mathrm{stop})=4
\]

\[
Q((4,C),\mathrm{draw})=3
\]

پس کنش حریصانه \(\mathrm{stop}\) است. با \(\epsilon=0.4\) و دو کنش ممکن داریم:

\[
P(\mathrm{stop})=(1-\epsilon)+\frac{\epsilon}{2}
=
0.6+0.2=0.8
\]

\[
P(\mathrm{draw})=\frac{\epsilon}{2}=0.2
\]

مقدار مورد انتظار \(Q\) بر اساس مقدارهای یادگرفته‌شده‌ی فعلی برابر است با:

\[
\mathbb{E}[Q]
=
0.8(4)+0.2(3)
=
3.8
\]

اگر مقدار واقعی \(Q^*\) را معیار بگیریم، چون در حالت \((4,C)\) با \lr{draw} حتماً می‌بازیم، داریم:

\[
\mathbb{E}[Q^*]
=
0.8(4)+0.2(0)
=
3.2
\]

\part[\textbf{و)}]
تعریف \lr{regret} هر گام به صورت زیر است:

\[
\mathrm{regret}_t=V^*(s_t)-Q^*(s_t,a_t)
\]

برای گام اول:

\[
s_1=(0,C),\qquad a_1=\mathrm{draw}
\]

\[
\mathrm{regret}_1
=
V^*(0,C)-Q^*((0,C),\mathrm{draw})
=
\frac{10}{3}-\frac{10}{3}
=
0
\]

برای گام دوم:

\[
s_2=(2,C),\qquad a_2=\mathrm{stop}
\]

\[
\mathrm{regret}_2
=
V^*(2,C)-Q^*((2,C),\mathrm{stop})
=
3-2
=
1
\]

برای گام سوم:

\[
s_3=(0,C),\qquad a_3=\mathrm{draw}
\]

\[
\mathrm{regret}_3
=
V^*(0,C)-Q^*((0,C),\mathrm{draw})
=
0
\]

برای گام چهارم:

\[
s_4=(4,C),\qquad a_4=\mathrm{draw}
\]

\[
\mathrm{regret}_4
=
V^*(4,C)-Q^*((4,C),\mathrm{draw})
=
4-0
=
4
\]

پس \lr{regret} تجمعی برابر است با:

\[
\mathrm{Regret}
=
0+1+0+4
=
5
\]

\[
\boxed{\mathrm{Regret}=5}
\]

\end{parts}
\clearpage

\section*{سؤال ششم}
\begin{parts}

\part[\textbf{الف)}]
برای به‌دست آوردن مدل گذار، تعداد رفتن‌ها از هر \((s,a)\) به حالت بعدی را حساب می‌کنیم و بعد تقسیم بر کل دفعات همان \((s,a)\) می‌کنیم:

\[
\hat{T}(s,a,s')=\frac{N(s,a,s')}{N(s,a)}
\]

برای حالت \(X\)، کنش \(a\) سه بار آمده است؛ دو بار به \(Y\) و یک بار به \(Z\):

\[
\hat{T}(X,a,Y)=\frac{2}{3},\qquad
\hat{T}(X,a,Z)=\frac{1}{3}
\]

برای کنش \(b\) در حالت \(X\)، هر دو بار به خود \(X\) رفته‌ایم:

\[
\hat{T}(X,b,X)=1
\]

برای حالت \(Y\)، کنش \(a\) سه بار آمده است؛ دو بار به \(Z\) و یک بار به \(Y\):

\[
\hat{T}(Y,a,Z)=\frac{2}{3},\qquad
\hat{T}(Y,a,Y)=\frac{1}{3}
\]

برای کنش \(b\) در حالت \(Y\)، فقط یک بار آمده و به \(X\) رفته است:

\[
\hat{T}(Y,b,X)=1
\]

برای حالت \(Z\)، کنش \(a\) دو بار آمده و هر دو بار در خود \(Z\) مانده است:

\[
\hat{T}(Z,a,Z)=1
\]

برای کنش \(b\) در حالت \(Z\)، فقط یک بار آمده و به \(X\) رفته است:

\[
\hat{T}(Z,b,X)=1
\]

بقیه‌ی حالت‌هایی که در داده نیامده‌اند، احتمالشان صفر در نظر گرفته می‌شود.

\part[\textbf{ب)}]
برای پاداش، برای هر حالت و کنش، میانگین پاداش‌های همان قسمت را می‌گیریم:

\[
\hat{R}(s,a,s')=\text{average reward for }(s,a,s')
\]

اینجا پاداش فقط به \(s\) و \(a\) بستگی دارد و حالت بعدی مقدارش را عوض نکرده است. پس داریم:

\[
\hat{R}(X,a,Y)=3,\qquad
\hat{R}(X,a,Z)=3
\]

\[
\hat{R}(X,b,X)=1
\]

\[
\hat{R}(Y,a,Z)=4,\qquad
\hat{R}(Y,a,Y)=4
\]

\[
\hat{R}(Y,b,X)=2
\]

\[
\hat{R}(Z,a,Z)=0
\]

\[
\hat{R}(Z,b,X)=5
\]

برای بقیه‌ی گذارها هم چون احتمالشان صفر است، لازم نیست پاداشی حساب کنیم.

\part[\textbf{ج)}]
فرمول یک مرحله از \lr{Value Iteration} این است:

\[
V_1(s)=\max_a \sum_{s'} \hat{T}(s,a,s')
\left[
\hat{R}(s,a,s')+\gamma V_0(s')
\right]
\]

چون در صورت سؤال گفته شده است که \(V_0(s)=0\)، قسمت \(\gamma V_0(s')\) صفر می‌شود. پس فقط میانگین پاداش هر کنش را حساب می‌کنیم:

\[
V_1(s)=\max_a \sum_{s'} \hat{T}(s,a,s')\hat{R}(s,a,s')
\]

برای حالت \(X\):

\[
Q_1(X,a)
=
\frac{2}{3}(3)+\frac{1}{3}(3)
=
3
\]

\[
Q_1(X,b)=1(1)=1
\]

پس:

\[
V_1(X)=\max\{3,1\}=3
\]

\[
\pi(X)=a
\]

برای حالت \(Y\):

\[
Q_1(Y,a)
=
\frac{2}{3}(4)+\frac{1}{3}(4)
=
4
\]

\[
Q_1(Y,b)=1(2)=2
\]

پس:

\[
V_1(Y)=\max\{4,2\}=4
\]

\[
\pi(Y)=a
\]

برای حالت \(Z\):

\[
Q_1(Z,a)=1(0)=0
\]

\[
Q_1(Z,b)=1(5)=5
\]

پس:

\[
V_1(Z)=\max\{0,5\}=5
\]

\[
\pi(Z)=b
\]

در نتیجه داریم:

\[
V_1(X)=3,\qquad \pi(X)=a
\]

\[
V_1(Y)=4,\qquad \pi(Y)=a
\]

\[
V_1(Z)=5,\qquad \pi(Z)=b
\]

\part[\textbf{د)}]
بله، \lr{Value Iteration} در این MDP همگرا می‌شود. دلیلش این است که تعداد حالت‌ها و کنش‌ها محدود است و مقدار ضریب تخفیف هم کمتر از ۱ است:

\[
\gamma=0.9<1
\]

یعنی اثر پاداش‌های خیلی دورتر، کم‌کم کمتر می‌شود و مقدارها به یک جواب ثابت می‌رسند. این‌که حالت \(Z\) با کنش \(a\) در خودش می‌ماند، مشکلی ایجاد نمی‌کند؛ چون باز هم ضریب تخفیف کمتر از ۱ است.

\end{parts}

% ---------------------------
\end{document}
